%! ~~~ Packages Setup ~~~ 
\documentclass[]{../../../../tex/classes/styledArticle}
\usepackage{../../../../tex/packages/hebrewSupport}
\usepackage{../../../../tex/packages/mathShortcuts}
\usepackage{../../../../tex/packages/theoremsSupport}
\usepackage{../../../../tex/packages/enfancysections}

%! ~~~ Document ~~~

\author{שחר פרץ}
\title{\textit{לינארית 2א $\sim$ חוגים ושאר ירקות}}
\date{24 לאפריל 2025}
\begin{document}
	\maketitle
	\section{\en{Rings}}
	\defi{תחום שלמות הוא חוג קומוטטיבי עם יחידה ללא מחלקי $0$. }
	\defi{חוק ייקרא \textit{ללא מחלקי 0} אם: \hfill $\forall a, b \in \R \co ab = 0 \implies a = 0 \lor b = 0$}
	
	דוגמאות לחוגים עם מחלקי $0$: 
	\begin{itemize}
		\item $M_2(\R)$: הוכחה $a = b = \binom{0\, 1}{0\, 0}, \ a \cdot b = 0$
		\item $\Z/_{6\Z}$ הוכחה $2 \cdot 3 = 0$. 
	\end{itemize}
	
	\theo{בתחום שלמות יש את כלל הצמצום בכפל: אם $ab = ac \land a \neq 0$ אז $b = c$. }\begin{proof}
		\[ ab  \cdot ac = 0 \implies a(b \cdot c) = 0 \implies a = 0 \lor b - c = 0 \]
		בגלל ש־$a \neq 0$, אז $b - c = 0$. נוסיף את $c$ הנגדי של $-c$ ונקבל $b = c$. 
	\end{proof}
	
	דוגמאות לתחום שלמות: 
	\begin{itemize}
		\item שדות
		\item השלמים
		\item חוג הפולינומים
	\end{itemize}
	
	\theo{לכל $f, g \in \F[x]$, אם $g \neq 0$ אז קיימים ויחידים פולינומים $q, r \n \F[x]$ כך ש־$f = qg + r \land \deg r < \deg g$}
	
	\defi{נאמר שפולינום $q$ מחלק את $f$ אם $r = 0$ ומסמנים $q \mid f$. }
	
	\defi{חוק אוקלידי הוא חוג שמעליו אפשר לבצע פירוק פולינום כזה. }
	
	דוגמה לחוג שאינו אוקלידי: $\Z[\sqrt{-5}]$ הוא $\{a + b\sqrt{-5} \mid a, b \in \Z\}$. 
		
	\theo{חוג אוקלידי $\impliedby$ פריקות יחידה (דומה למשפט היסודי של האריתמטיקה). }
		
		לדוגמה בחוג לעיל $6 = 2 \cdot 3 = (1 + \sqrt{-5})(1 - \sqrt{-5})$ על אף ש־$2, 3$ אי־פריקים וכן $(1 + \sqrt{-5}), (1  - \sqrt{-5})$ אי פריקים. 
		
	\cola{\,
	\begin{itemize}
		\item $f(a) = 0 \iff (x - a) \mid f$ (משפט בזו)
		\item אם $\deg f = n > -\inf$, ל־$f$ לכל היותר $n$ שורשים כולל ריבוי. 
		\item נניח ש־$f, g \in \F[x]$ ו־$F \subseteq K$, כאשר $K$ שדה. אם $g \mid f$ מעל $K$ אז $g \mid f$ מעל $\F$. 
	\end{itemize}
	}
	\begin{proof}
		\begin{enumerate}\,
			\item \begin{itemize}
				\item[$\implies$] נניח $x - a \mid f$. אז קיים פולינום $g$ כך ש־$f = (x - a)g$. אז $f(a) = (a - a)g(a) = 0$. 
				\item[$\impliedby$] נניח $f(a) = 0$. אז קיימים $q, r \in \F[x]$ כך ש־$f = q(x - a)$ ועל כן $0 = f(a) = q(a)(a - a) + r(a) = 0$ ולכן $r(a) = 0$. משום ש־$r$ פולינום קבוע (דרגתו קטנה מ־1, כי חילקנו ב־$(x - a)$ מדרגה 1), אז $r(x) = 0$. 
			\end{itemize}
			\item אינדוקציה
			\item נוכיח ב"contrapositive": אנו יודעים ש־$P \to Q \iff \lnot Q \to \lnot P$. נניח ש־$g \nmid f$ מעל $\F$. קיימים $q, r \in \F[x]$ כך ש־$f =qg + r, \ r \neq 0$. הפירוק הזה הוא גם ב־$K[x]$. מיחידות $r$, נקבל ש־$g \nmid f$ כל מעל $K$. 
		\end{enumerate}
	\end{proof}
	
	\subsection{עוד על תחומי שלמות}
	\defi{יהי $R$ תחום שלמות, $a, b \in R$. נאמר ש־$a \mid b$ אם קיים $c \in \R$ כך ש־$ac = b$. }
	\defi{$u \in R$ נקרא \textit{הפיך} אם קיים $\ag \in R$ כך ש־$\ag u = 1$. }
	\theo{יהי $R$ תחום שלמות, $u \in R$ הפיך. יהי $a \in \R$. אז $u \mid a$. }\begin{proof}
		$1 \mid a, \ u \mid 1$. יחס החלוקה טרנזטיבי ולכן $u \mid a$. 
	\end{proof}
	\noti{קבוצת ההפיכים מוסמנת ב־$R^x$. }
	\textbf{דוגמאות. }
	\begin{enumerate}
		\item אם $R = \F$, אז $\F^x = \F \setminus \{0\}$
		\item אם $R = \Z$ אז $\Z^2 = \{\pm 1\}$
		\item אם $R = \F[x]$ אז $R^x = \F^x$ (ההתייחסות לסקלרים $\F$ היא כאל פונקציות קבועות)
	\end{enumerate}
	\defi{$a, b \in R$ נקראים \textit{חברים} אם קיים $u \in R^x$ הפיך כך ש־$a = ub$, ומסמנים $a \sim b$ }
	"אני אני חבר של עומר, ועומר חבר של מישהו בכיתה שאני לא מכיר, אני לא חבר של מי שאני לא מכיר." המרצה: "למה לא? תהיה חבר שלו". 
	
	\theo{יחס החברות הוא יחס שקילות. }\begin{proof}
		\begin{enumerate}[A.]
			\item $a \sim a$ כי $1 \in R^x$
			\item אם $a \sim b$ אז קיים $u \in R^x$ כך ש־$a = ub$. קיים ל־$u$ הופכי $\ag$ אז $\ag a + \ag u b = b$ ולכן $b \sim a$. 
			\item נניח $a \sim b \land b \sim c$, כי מכפלת ההופכיים הפיכה $a \sim c$ וסיימנו. 
		\end{enumerate}
	\end{proof}
	\theo{הופכי הוא יחיד} (אותה ההוכחה כמו בשדה. לא בהכרח בתחום שלמות, מעל כל חוג)\begin{proof}
				יהי $a \in R^x$ ו־$u, u'$ הופכיים שלו, אז:
		\[ u = u \cdot 1 = u \cdot a \cdot u' = 1 \cdot u' = u' \]
	\end{proof}
	\theo{אם $a \mid b$ וכם $b \mid a$ אז $a \mid b$ (בתחום שלמות). }\begin{proof}
		\begin{align*}
			a \mid b &\implies \exists c \in \R \co ac = b \\
			b \mid a &\implies \exists d \in \R\co bd = a
		\end{align*}
		לכן: 
		\[ ac = b \implies acd = a \implies a(cd - 1) = 0 \implies a = 0 \lor cd = 1 \]
	\end{proof}
	אם $a = 0$ אז $b = 0$ (ממש לפי הגדרה) ו־$\sim$ שקילות (רפליקסיביות). אחרת, $cd = 1$ ולכן $c$ הפיך, סה"כ $a \mid b$. 
	
	"אני חושב שבעברית קראו להם ידידים, לא רצו להתחייב לחברות ממש". 
	
	
	\defi{איבר $p \in R$ נקרא \textit{אי־פריק} אם מתקיים $p = ab \implies a \in R^x \lor b \in R^x$. }
	
	\defi{איבר $p \in R$ יקרא \textit{ראשוני} אם $p \mid (a \cdot b) \implies p \mid a \lor p \mid b$. }
	\textit{הערה: }איברים הפיכים לא נחשבים אי־פריקים או ראשוניים. הסיבה להגדרה: בשביל נכונות המשפט היסודי של האריתמטיקה (יחידות הפירוק לראשוניים). 
	
	\theo{בתחום שלמות כל ראשוני הוא אי פריק. }
	\textit{הערה: }שקילות לאו דווקא. 
	\begin{proof}
		יהי $p \in R$ ראשוני. יהיו $a, b \in R$ כך ש־$p = ab$. בה"כ $p \mid a$. אז קיים $c$ כך ש־$pc = a$ ולכן $pcb = p$. סה"כ $p \neq 0$ ולכן $cb = 1$ (ראה לעיל) ו־$b$ הפיך. 
	\end{proof}
	\theo{נניח שבתחום שלמות $R$, כל אי־פריק הוא גם ראשוני. אז $R$ תחום פריקות יחידה. }
	\defi{$R$ תחום פירוקת יחידה אם $\prod_{i = 1}^{n}p_i = \prod_{j = 1}^{m}q_i$ עבור $p_i, q_j$ ראשוניים, אז $m = n$, ועד לכדי סידור מחדש, לכל $i \in [n]$ $p_i \sim q_i$}
	
	ההוכחה: זהה לחלוטין לזו של המשפט היסודי. \begin{proof}
		באינדוקציה על $n+m$. בסיס: $n + m = 2$ ולכן $n = m = 1$ (כי מעפלה ריקה לא רלוונטית מאוד) אז $p = q$. נעבור לצעד. נניח שהטענה נכונה לכל $n + m < k$. נניח ש־$n + m = k$. אז $p_1 \mid \prod_{j = 1}^{m}q_j$. בה"כ $p_1 \mid q_1$. $q_1$ אי־פריק ולא הפיך. $p_1$ לא הפיך. לכן $p_1 \sim q_1$. אז עד כדי כפל בהופכי נקבל ש־$\prod_{i = 2}^{n} p_i = \prod_{j = 2}^{n} q_j$. \textit{הערה: }ראשוני כפול הפיך נשאר ראשוני. מכאן הקענו לדרוש וסיימנו (\textit{הערה שלי: }כאילו תכפילו בחברים ותקבלו את מה שצריך). 
	\end{proof}
	
	\defi{יהי $R$ תחום שלמות. תת־קבוצה $0 \neq I \subseteq R$ נקראת \textit{אידיאל} אם: 
	\begin{itemize}[A.]
		\item $\forall a, b \in I \co a + b \in I$ – סגירות לחיבור. 
		\item $\forall a \in I \, \forall b \in R \co ab \in I$ – תכונת הבליעה. [בפרט $0 \in I$]
	\end{itemize}}
	\textbf{דוגמאות: }\begin{enumerate}
		\item $0$ תמיד אידיאל, כך החוג תמדי אידיאל. 
		\item הזוגיים ב־$\Z$. 
		\item לכל $n \in \Z$, $n\Z$ אידיאל ($n$ כפול השלמים). הזוגיים מקרה פרטי. 
		\item $\la f \ra \subseteq \F[x]$ המוגדר לפי $\la f \ra := \{g \in \F[x] \mid f|g\}$
		\item הכללה של הקודמים: עבור $a \in R$ נסמן $\la a \ra := \{a \cdot b \mid b \in R\}$
		\item $I = \{f \in \F[x] \mid f(0) = 0\}$ (לעיתים מסומן $\forall a \in R \co aR = \la a \ra$)
		\item נוכל להכליל את 4 עוד: ("הכללה של הכללה היא הכללה. זה סגור להכללה. זה קורה הרבה במתמטיקה")
		\[ I = aR + bR = \{ar + bs \mid r, s \in R\} \]
		וניתן להכליל עוד באינדוקציה. 
	\end{enumerate}
	\defi{אידיאל $I$ נקרא \textit{ראשי} אם הוא מהצורה $aR$ עבור $a \in R$ כלשהו. }
	\defi{תחום שלמות נקרא \textit{ראשי} אם כל אידיאל שלו ראשי. }
	\theo{נניח ש־$R$ ראשי, אז כל אי־פריק ב־$R$ הוא ראשוני. }
	(תנאי מספיק אך לא הכרחי)
	
	\begin{proof}
		יהי $p \in R$ אי פריק. נראה שהוא ראשוני. נביט ב־$ab \in R$ כך ש־$p \mid ab$. נביט ב־$I = aR + bR$. מהיותו ראשי קיים איזשהו $c \in R$ כך ש־$I = cR$. אז $c \in I$ (כי $c \cdot 1 \in R$). אז $a, p \in I$. $a$ ברור. קיים $d \in R$ כך ש־$pd = a \cdot b$. למה $p \in I$ – המרצה דפק קליגמן ולא יודע להוכיח. 
		אחרי משהו כמו 5 דק' של בהיה בלוח הוא בה עם הדבר הבא: 
		\[ c \in I \implies \exists r, s \in \R\co ar + bs = c \]
		"אני אחושב על זה ואני אמשיך פעם הבאה". 
	\end{proof}
	
	
	
	\ndoc
\end{document}
